% ==========================================================================
% GENERATED by scripts/split_paper.py. Do not edit this file.
%
% Source: imported_papers/sds_hounsfield_features/main.tex
% Part 6: Ablation Studies
%
% Edits belong upstream, in the published paper. To change what the
% chapter includes, change imports/hounsfield.toml and re-run
% the splitter.
% ==========================================================================
\section{Ablation Studies}

\subsection{Removing raw HU values}

To evaluate the sufficiency of the derived features $\bm{X}_c$ alone, we conduct experiments excluding the original HU values $\bm{X}_0$ (the original CT image) from the input. This modification reduces the available input information, as the binary masks are \emph{derived from} the HU values but are \emph{not invertible}.
However, as shown in \cref{tbl:results-hut}, the performance remains high across all model sizes, even surpassing the results obtained with raw HU values alone (\cref{tbl:results}).
This confirms that the inductive bias introduced by our ``feature engineering'' provides a more effective representation of the information needed for semantic segmentation of CT images with respect to BCA.

\begin{table}[t]
    \caption{Segmentation performance when trained using only binary masks (excluding raw HU values from the CT images).}
     \begin{center}
    \begin{tabular}{lrr}
        \toprule
        \multicolumn{1}{c}{\textbf{Model}} 
        & \multicolumn{1}{c}{\textbf{Dice score}} 
        & \multicolumn{1}{c}{\textbf{F$_1$-Score}} \\
        \midrule
        U-Net-small, (only) HU ranges                      
        % & \formatpercentage{0.8906}\ \ $\pm$\formatpercentage{0.0036} 
        & \formatpercentage{0.9424}\ \ $\pm$\formatpercentage{0.002643} 
        & \formatpercentage{0.9383}\ \ $\pm$\formatpercentage{0.003656}  \\
        \midrule
        U-Net-base,\ \ (only) HU ranges
        % & \formatpercentage{0.8933}\ \ $\pm$\formatpercentage{0.0021} 
        & \formatpercentage{0.9439}\ \ $\pm$\formatpercentage{0.002513} 
        & \formatpercentage{0.94}\ \ $\pm$\formatpercentage{0.002301} \\
        \midrule
        U-Net-large, (only) HU ranges                       
        % & \formatpercentage{0.8854}\ \ $\pm$\formatpercentage{0.0041}
        & \formatpercentage{0.9403}\ \ $\pm$\formatpercentage{0.003839}  
        & \formatpercentage{0.9356}\ \ $\pm$\formatpercentage{0.003553} \\
        \bottomrule
    \end{tabular}
    \label{tab:results-no-hu-values} \label{tbl:results-hut}
     \end{center}
\end{table}

\subsection{Removing individual binary masks during inference} \label{sec:remove-masks}

To further evaluate the contribution of individual binary masks, we measure the performance drop when artificially removing a single binary mask \emph{during inference} by setting all locations to zero. 
\Cref{tbl:results-removed-mask} presents the results using the final checkpoint of our \emph{U-Net-base, HU + HU ranges} model.

The middle column shows the drop in performance for a specific tissue type if its binary mask is removed.
The segmentation performance for SM and VAT drops substantially to \SI{0}{\percent} when we remove their corresponding binary mask at inference, suggesting that the trained model relies heavily on the binary masks to segment these tissue types.
Since IMAT and SAT share identical thresholds (\Cref{tab:hu_values}), we remove their (identical) binary masks together.
Removing those masks results in a notable decrease in segmentation performance, particularly for IMAT, where the score drops by over \SI{70}{\percent}.
Overall, the results in the middle column suggest that the model utilizes a tissue's binary mask to segment the specific tissue type.

The right column shows the drop in performance for the \emph{remaining} tissue classes (all tissue classes whose masks we did not remove).
Notably, this score is also affected, dropping \SI{24}{\percent} if we remove the SM binary mask and around \SIrange{8}{9}{\percent} if removing other masks.
While we suspect that the binary masks help the model to better detect other tissue types as well,
this overall performance drop on other tissue types could also be due to having unrealistic input, as the binary masks don't match the CT input image, and thus are out of distribution.
% introduced by our feature engineering is not only leveraged by the model but also contributes strongly to the model's final performance.

\begin{table}[h]
    \caption{
    Effect of removing individual binary masks during inference. Numbers (left $\rightarrow$ right) show the Dice score drop from the unchanged binary mask (left) to the removed binary mask (right).
    The first column shows which tissue's binary mask we remove.
    The second column shows the Dice score drop for the tissue classes whose mask we remove.
    The third column shows the Dice score drop for the \emph{remaining} tissue classes (all tissue classes whose masks we did not remove).
    Last row: As IMAT and SAT share the same HU range, they are removed together, with the first tissue Dice score for IMAT and the second for SAT.
    }
    \begin{center}
    \begin{tabular}{lrr}
        \toprule
        \multicolumn{1}{c}{\textbf{Removed mask}} 
        & \multicolumn{1}{c}{\textbf{Tissue Dice score}}
        & \multicolumn{1}{c}{\textbf{Other Dice score}} \\
        \midrule
        SM
        & \formatpercentage{0.9508} $\rightarrow$\ \ \ \SI{0.0}{\percent}     
        & \formatpercentage{0.9338} $\rightarrow$ \formatpercentage{0.6930}  \\   
        \midrule
        VAT
        & \formatpercentage{0.9409} $\rightarrow$\ \ \ \SI{0.0}{\percent}
        & \formatpercentage{0.9344} $\rightarrow$ \formatpercentage{0.8476} \\

        \midrule
        \multirow{2}{*}{\shortstack[l]{IMAT and \\ SAT}}
        %\midrule
        %IMAT (+ SAT)
        & \formatpercentage{0.8322} $\rightarrow$ \formatpercentage{0.1064}
        & \multirow{2}{*}{\formatpercentage{0.9484} $\rightarrow$ \formatpercentage{0.8651}} \\
        & \formatpercentage{0.9800} $\rightarrow$ \formatpercentage{0.6846} & \\
        \bottomrule
    \end{tabular}
    \label{tbl:results-removed-mask}
    \end{center}
\end{table}
